Privacy parameters and watermark decisions do not establish assurance when their
targets, mechanisms, interfaces, and evidence are misaligned.  This dissertation
develops Claim-Aligned Assurance through Auditable Privacy and Secure Public
Verification.

The Auditable Privacy study takes frozen evidence from a completed windowed
DP-SGD execution together with requested wording and decides which privacy-unit
statement the execution can license.  It retains an aligned native claim, applies
a registered conversion or group consequence when its premises are complete, and
blocks positive wording when evidence is missing or the numerical consequence is
vacuous.  The complete validator matched the expected tuple on 100 allowed and
100 paired non-allow cases.  Bound WISDM and UCI HAR executions demonstrated
query-dependent authorization, while mapping diagnostics over 5,000 Sepsis
patients and more than 27.7 million Backblaze drive-day records exposed how window
policy changes raw-to-generated influence. 

PP-Mark combines a context-bound latent watermark, trace commitment, and a
zero-knowledge proof receipt.  Full acceptance requires a signal and receipt for
the same context and image.  A 10,000-binding search and white-box optimization
exposed score attacks, but no full acceptance was observed; transformation
true-positive rates were 0.74--1.00 in score mode.

Across the two studies, the dissertation derives claim specification, mechanism
binding, evidence-bound decision, and fail-closed assurance.  It limits positive
privacy and provenance statements to the objects, mechanisms, interfaces, and
evidence that jointly support them.
